\section{Dynamic regularity of the process}\label{sec:regularidad}

The entropy comparison already allows us to transfer sufficiently rare events. To apply it to repairs, we
first need to know how often
the process selects the lines that the two queens in a repair would occupy. These probabilities depend on
the weights and the available
degrees at each step.

\subsection{Rescaled weights and initial balance}

For $1\le h\le K$ we rescale only the real weights by the common factor $x_*^{-2}$, which leaves the law
unchanged, giving
$v_e=\varpi(e)/x_*^2=r_Rr_Cu_Du_A$. For $h=0$, $v_e$ is the real weight from
Definition~\ref{def:referencia_cero} itself. By
Theorem~\ref{thm:perfil}(a), $c\le v_e\le C$ with $c=\underline r^2\kappa^2$ and $C=M^2$. Set \[
 W_0=\sum_ev_e\asymp n^2,\quad \deg_v^0=\sum_{e\ni v}v_e,\quad \pi_e=\frac{qv_e}{W_0},\quad
 p_v=\sum_{e\ni v}\pi_e,\quad U_f^0=\sum_{v\in f}u_v .
\] We retain the virtual weights $u_v$. The dummy vertices have $u_v=1$ and $p_v=0$.

\begin{lema}[Initial balance]\label{lem:balance} Uniformly in $0\le h\le K$, for every physical vertex $v$
and every real edge $e$,
\begin{equation}\label{eq:balance}
 p_v+u_v=1+O(n^{-1}),\qquad U_f^0=H+O(1),\qquad p_v\le C,\qquad \pi_e\le C/n .
\end{equation} Moreover, $\sum_{v\in f}p_v=q$ in each family, and for diagonals and antidiagonals,
$p_v\le1-\eta$ with $\eta=\kappa/2$,
for $n$ large. \end{lema}

\begin{proof} Let $\rho_n=q/n$. By quadrature, $W_0=n^2\rho_n/b\,(1+O(n^{-1}))$. On a row $r$, the Euler
equation~\eqref{eq:euler} gives
$q\deg_v^0/W_0=b\,r_R(\bar r)J_R(\bar r)+O(n^{-1})=1-u_R+O(n^{-1})$. On a diagonal,
$q\deg_v^0/W_0=b\,u_DJ_D+O(n^{-1})=1-u_D+O(n^{-1})$. The arguments for columns and antidiagonals are
analogous. The bounds and
quadrature estimates are uniform up to the boundary of the target profile. Each real edge appears exactly
once in each family, so
$\sum_{v\in f}\deg_v^0=W_0$ and $\sum_{v\in f}p_v=q$. Summing the pointwise balance gives
$U_f^0=N-q+O(N/n)=H+O(1)$, including the
dummy vertices. On diagonals, $u_v\ge\kappa$, whence $p_v\le1-\kappa+O(n^{-1})\le1-\eta$. At $h=0$ the
balance is $p_v=1+O(n^{-1})$,
with $u_v=0$ on physical rows and columns. The sum of the virtual weights in each of these families is $H$
(from dummy vertices alone),
and the diagonals are unchanged. \end{proof}

\emph{Clocks.} The index $t$ counts all pieces placed, $z_t=(N-t)/N$ is the active fraction, and $q_t$ is
the number of real pieces still to be
placed. Neither $z_t$ nor $q_t/n$ replaces the fixed parameter $x_*$ in the weights.
